\section{T}

% tauschen (to exchange / swap)
\begin{verbentry}{
  style = {default},
  toc = {tauschen (to exchange / swap)},
  name = {tauschen},
  english = {to exchange / swap},
  type = {regular},
  auxiliary = {haben}
}


\vspace{5pt}
\hrule
\vspace{5pt}

\begin{tabular}{rl}
\multicolumn{2}{l}{\textbf{PRÄSENS}}\\
ich & \textbf{tausche}\\
du & \textbf{tauschst}\\
er/sie/es & \textbf{tauscht}\\
wir & \textbf{tauschen}\\
ihr & \textbf{tauscht}\\
sie/Sie & \textbf{tauschen}\\
\end{tabular}
\hfill
\begin{tabular}{rl}
\multicolumn{2}{l}{\textbf{PRÄTERITUM}}\\
ich & \textbf{tauschte}\\
du & \textbf{tauschtest}\\
er/sie/es & \textbf{tauschte}\\
wir & \textbf{tauschten}\\
ihr & \textbf{tauschtet}\\
sie/Sie & \textbf{tauschten}\\
\end{tabular}
\hfill
\begin{tabular}{rl}
\multicolumn{2}{l}{\textbf{IMPERATIV}}\\
& \textbf{tausche (du)!}\\
& \textbf{tauschen wir!}\\
& \textbf{tauscht (ihr)!}\\
& \textbf{tauschen Sie!}\\
\end{tabular}

\vspace{10pt}

\begin{tabular}{rl}
\multicolumn{2}{l}{\textbf{PERFEKT}}\\
& haben + \textbf{getauscht}\\
\end{tabular}
\hfill
\begin{tabular}{rl}
\multicolumn{2}{l}{\textbf{FUTURE I}}\\
& werden + \textbf{tauschen}\\
\end{tabular}
\hfill
\begin{tabular}{rl}
\multicolumn{2}{l}{\textbf{PLUSQUAMPERFEKT}}\\
& hatten + \textbf{getauscht}\\
\end{tabular}

\vspace{10pt}

\begin{tabular}{lll}
\multicolumn{3}{l}{\textbf{MIT PRÄPOSITIONEN}}\\
& \textbf{tauschen gegen} + Akk & (to exchange for)\\
& \textbf{tauschen mit} + Dat & (to exchange with)\\
\end{tabular}
\hfill
\begin{tabular}{lll}
\multicolumn{3}{l}{\textbf{TRENNBARE VERBEN}}\\
& \textbf{aus$|$tauschen} & (to exchange / replace)\\
& \textbf{um$|$tauschen} & (to return (a purchased item))\\
\end{tabular}
\vspace{-5pt}
\end{verbentry}

% teilen (to share / divide)
\begin{verbentry}{
  style = {default},
  toc = {teilen (to share / divide)},
  name = {teilen},
  english = {to share / divide},
  type = {regular},
  auxiliary = {haben}
}


\vspace{5pt}
\hrule
\vspace{5pt}

\begin{tabular}{rl}
\multicolumn{2}{l}{\textbf{PRÄSENS}}\\
ich & \textbf{teile}\\
du & \textbf{teilst}\\
er/sie/es & \textbf{teilt}\\
wir & \textbf{teilen}\\
ihr & \textbf{teilt}\\
sie/Sie & \textbf{teilen}\\
\end{tabular}
\hfill
\begin{tabular}{rl}
\multicolumn{2}{l}{\textbf{PRÄTERITUM}}\\
ich & \textbf{teilte}\\
du & \textbf{teiltest}\\
er/sie/es & \textbf{teilte}\\
wir & \textbf{teilten}\\
ihr & \textbf{teiltet}\\
sie/Sie & \textbf{teilten}\\
\end{tabular}
\hfill
\begin{tabular}{rl}
\multicolumn{2}{l}{\textbf{IMPERATIV}}\\
& \textbf{teil(e) (du)!}\\
& \textbf{teilen wir!}\\
& \textbf{teilt (ihr)!}\\
& \textbf{teilen Sie!}\\
\end{tabular}

\vspace{10pt}

\begin{tabular}{rl}
\multicolumn{2}{l}{\textbf{PERFEKT}}\\
& haben + \textbf{geteilt}\\
\end{tabular}
\hfill
\begin{tabular}{rl}
\multicolumn{2}{l}{\textbf{FUTURE I}}\\
& werden + \textbf{teilen}\\
\end{tabular}
\hfill
\begin{tabular}{rl}
\multicolumn{2}{l}{\textbf{PLUSQUAMPERFEKT}}\\
& hatten + \textbf{geteilt}\\
\end{tabular}

\vspace{10pt}

\begin{tabular}{lll}
\multicolumn{3}{l}{\textbf{MIT PRÄPOSITIONEN}}\\
& \textbf{teilen mit} + Dat & (to share with)\\
& \textbf{teilen in} + Akk & (to divide into)\\
\end{tabular}
\hfill
\begin{tabular}{lll}
\multicolumn{3}{l}{\textbf{TRENNBARE VERBEN}}\\
& \textbf{auf$|$teilen} & (to split up / divide)\\
& \textbf{ab$|$teilen} & (to allocate / section off)\\
\end{tabular}
\vspace{-5pt}
\end{verbentry}

% tragen (to wear, carry)
\begin{verbentry}{
  style = {acc},
  toc = {tragen (to wear, carry)},
  name = {tragen},
  english = {to wear, carry},
  type = {irregular, + Akkusativ},
  auxiliary = {haben}
}
 
    
     \vspace{5pt}

    \hrule
    
    \vspace{5pt}

  \begin{tabular}{rl}
     \multicolumn{2}{l}{\textbf{PRÄSENS} }\\
    ich & \textbf{trage}\\
    du & \textbf{trägst}\\
    er/sie/es & \textbf{trägt}\\
    wir & \textbf{tragen}\\
    ihr & \textbf{tragt}\\
    sie/Sie & \textbf{tragen}\\
  \end{tabular}
  \hfill
     \begin{tabular}{rl}
    \multicolumn{2}{l}{\textbf{PRÄTERITUM} }\\
    ich & \textbf{trug}\\
    du & \textbf{trugst}\\
    er/sie/es & \textbf{trug}\\
    wir & \textbf{trugen}\\
    ihr & \textbf{trugt}\\
    sie/Sie & \textbf{trugen}\\
  \end{tabular}
  \hfill
  \begin{tabular}{rl}
     \multicolumn{2}{l}{\textbf{IMPERATIV} }\\
   & \textbf{trag(e) (du)!}\\
   & \textbf{tragen wir!}\\
   & \textbf{tragt (ihr)!}\\
   & \textbf{tragen Sie!}\\
   & \vspace{10pt}
  \end{tabular}

    \vspace{10pt}

 \begin{tabular}{rl}
     \multicolumn{2}{l}{\textbf{PERFEKT}}\\
   & haben + \textbf{getragen}\\
  \end{tabular}
\hfill
   \begin{tabular}{rl}
   
    \multicolumn{2}{l}{\textbf{FUTURE I} }\\
   & werden + \textbf{tragen}\\
 
  \end{tabular}
\hfill
   \begin{tabular}{rl}
      \multicolumn{2}{l}{\textbf{PLUSQUAMPERFEKT}}\\
   & hatten + \textbf{getragen}\\
  \end{tabular}
  
  
\vspace{10pt}

\begin{tabular}{lll}
\multicolumn{3}{l}{\textbf{MIT PRÄPOSITIONEN}}\\
& \textbf{---} & \\
\end{tabular}
\hfill
\begin{tabular}{lll}
\multicolumn{3}{l}{\textbf{TRENNBARE VERBEN}}\\
& \textbf{---} & \\
\end{tabular}
\vspace{-5pt}
\end{verbentry}

% trainieren (to train / practice)
\begin{verbentry}{
  style = {default},
  toc = {trainieren (to train / practice)},
  name = {trainieren},
  english = {to train / practice},
  type = {regular},
  auxiliary = {haben}
}
 
    
     \vspace{5pt}
    \hrule
    \vspace{5pt}

  \begin{tabular}{rl}
     \multicolumn{2}{l}{\textbf{PRÄSENS} }\\
    ich & \textbf{trainiere}\\
    du & \textbf{trainierst}\\
    er/sie/es & \textbf{trainiert}\\
    wir & \textbf{trainieren}\\
    ihr & \textbf{trainiert}\\
    sie/Sie & \textbf{trainieren}\\
  \end{tabular}
  \hfill
     \begin{tabular}{rl}
    \multicolumn{2}{l}{\textbf{PRÄTERITUM} }\\
    ich & \textbf{trainierte}\\
    du & \textbf{trainiertest}\\
    er/sie/es & \textbf{trainierte}\\
    wir & \textbf{trainierten}\\
    ihr & \textbf{trainiertet}\\
    sie/Sie & \textbf{trainierten}\\
  \end{tabular}
  \hfill
  \begin{tabular}{rl}
     \multicolumn{2}{l}{\textbf{IMPERATIV} }\\
   & \textbf{trainier(e) (du)!}\\
   & \textbf{trainieren wir!}\\
   & \textbf{trainiert (ihr)!}\\
   & \textbf{trainieren Sie!}\\
   & \vspace{10pt}
  \end{tabular}

 \vspace{10pt}

 \begin{tabular}{rl}
     \multicolumn{2}{l}{\textbf{PERFEKT}}\\
   & haben + \textbf{trainiert}\\
  \end{tabular}
\hfill
   \begin{tabular}{rl}
    \multicolumn{2}{l}{\textbf{FUTURE I}}\\
   & werden + \textbf{trainieren}\\
  \end{tabular}
\hfill
   \begin{tabular}{rl}
      \multicolumn{2}{l}{\textbf{PLUSQUAMPERFEKT}}\\
   & hatten + \textbf{trainiert}\\
  \end{tabular}

   \vspace{10pt}

   \begin{tabular}{lll}
      \multicolumn{3}{l}{\textbf{MIT PRÄPOSITIONEN}}\\
& \textbf{trainieren für} + Akk & (to train for)\\
& \textbf{trainieren mit} + Dat & (to train with)\\
& \textbf{trainieren gegen} + Akk & (to train against / play against)\\
\end{tabular}
  \hfill
   \begin{tabular}{lll}
      \multicolumn{3}{l}{\textbf{TRENNBARE VERBEN}}\\
 & \textbf{ab$|$trainieren} & (to train off / reduce through training)\\
\vspace{20pt}
  \end{tabular}
  \vspace{-5pt}
\end{verbentry}

% träumen (to dream)
\begin{verbentry}{
  style = {default},
  toc = {träumen (to dream)},
  name = {träumen},
  english = {to dream},
  type = {regular},
  auxiliary = {haben}
}
 
    
     \vspace{5pt}

    \hrule
    
    \vspace{5pt}

  \begin{tabular}{rl}
     \multicolumn{2}{l}{\textbf{PRÄSENS} }\\
    ich & \textbf{träume}\\
    du & \textbf{träumst}\\
    er/sie/es & \textbf{träumt}\\
    wir & \textbf{träumen}\\
    ihr & \textbf{träumt}\\
    sie/Sie & \textbf{träumen}\\
  \end{tabular}
  \hfill
     \begin{tabular}{rl}
    \multicolumn{2}{l}{\textbf{PRÄTERITUM} }\\
    ich & \textbf{träumte}\\
    du & \textbf{träumtest}\\
    er/sie/es & \textbf{träumte}\\
    wir & \textbf{träumten}\\
    ihr & \textbf{träumtet}\\
    sie/Sie & \textbf{träumten}\\
  \end{tabular}
  \hfill
  \begin{tabular}{rl}
     \multicolumn{2}{l}{\textbf{IMPERATIV} }\\
   & \textbf{träum(e) (du)!}\\
   & \textbf{träumen wir!}\\
   & \textbf{träumt (ihr)!}\\
   & \textbf{träumen Sie!}\\
   & \vspace{10pt}
  \end{tabular}

    \vspace{10pt}

 \begin{tabular}{rl}
     \multicolumn{2}{l}{\textbf{PERFEKT}}\\
   & haben + \textbf{geträumt}\\
  \end{tabular}
\hfill
   \begin{tabular}{rl}
   
    \multicolumn{2}{l}{\textbf{FUTURE I} }\\
   & werden + \textbf{träumen}\\
 
  \end{tabular}
\hfill
   \begin{tabular}{rl}
      \multicolumn{2}{l}{\textbf{PLUSQUAMPERFEKT}}\\
   & hatten + \textbf{geträumt}\\
  \end{tabular}
  
  \vspace{10pt}


\begin{tabular}{lll}
\multicolumn{3}{l}{\textbf{MIT PRÄPOSITIONEN}}\\
& \textbf{---} & \\
\end{tabular}
\hfill
\begin{tabular}{lll}
\multicolumn{3}{l}{\textbf{TRENNBARE VERBEN}}\\
& \textbf{---} & \\
\end{tabular}

  \vspace{-5pt}
\end{verbentry}

% treten (to step)
\begin{verbentry}{
  style = {default},
  toc = {treten (to step)},
  name = {treten},
  english = {to step},
  type = {irregular},
  auxiliary = {sein}
}
 
    
     \vspace{5pt}

    \hrule
    
    \vspace{5pt}

  \begin{tabular}{rl}
     \multicolumn{2}{l}{\textbf{PRÄSENS} }\\
    ich & \textbf{trete}\\
    du & \textbf{trittst}\\
    er/sie/es & \textbf{tritt}\\
    wir & \textbf{treten}\\
    ihr & \textbf{tretet}\\
    sie/Sie & \textbf{treten}\\
  \end{tabular}
  \hfill
     \begin{tabular}{rl}
    \multicolumn{2}{l}{\textbf{PRÄTERITUM} }\\
    ich & \textbf{trat}\\
    du & \textbf{tratst}\\
    er/sie/es & \textbf{trat}\\
    wir & \textbf{traten}\\
    ihr & \textbf{tratet}\\
    sie/Sie & \textbf{traten}\\
  \end{tabular}
  \hfill
  \begin{tabular}{rl}
     \multicolumn{2}{l}{\textbf{IMPERATIV} }\\
   & \textbf{tritt (du)!}\\
   & \textbf{treten wir!}\\
   & \textbf{tretet (ihr)!}\\
   & \textbf{treten Sie!}\\
   & \vspace{10pt}
  \end{tabular}

    \vspace{10pt}

 \begin{tabular}{rl}
     \multicolumn{2}{l}{\textbf{PERFEKT}}\\
   & sein + \textbf{getreten}\\
  \end{tabular}
\hfill
   \begin{tabular}{rl}
   
    \multicolumn{2}{l}{\textbf{FUTURE I} }\\
   & werden + \textbf{treten}\\
 
  \end{tabular}
\hfill
   \begin{tabular}{rl}
      \multicolumn{2}{l}{\textbf{PLUSQUAMPERFEKT}}\\
   & waren + \textbf{getreten}\\
  \end{tabular}
  
  \vspace{10pt}



\begin{tabular}{lll}
\multicolumn{3}{l}{\textbf{MIT PRÄPOSITIONEN}}\\
& \textbf{---} & \\
\end{tabular}
\hfill
\begin{tabular}{lll}
\multicolumn{3}{l}{\textbf{TRENNBARE VERBEN}}\\
& \textbf{---} & \\
\end{tabular}

  \vspace{-5pt}
\end{verbentry}

% trinken (to drink)
\begin{verbentry}{
  style = {acc},
  toc = {trinken (to drink)},
  name = {trinken},
  english = {to drink},
  type = {irregular, + Akkusativ},
  auxiliary = {haben}
}
 
    
     \vspace{5pt}

    \hrule
    
    \vspace{5pt}

  \begin{tabular}{rl}
     \multicolumn{2}{l}{\textbf{PRÄSENS} }\\
    ich & \textbf{trinke}\\
    du & \textbf{trinkst}\\
    er/sie/es & \textbf{trinkt}\\
    wir & \textbf{trinken}\\
    ihr & \textbf{trinkt}\\
    sie/Sie & \textbf{trinken}\\
  \end{tabular}
  \hfill
     \begin{tabular}{rl}
    \multicolumn{2}{l}{\textbf{PRÄTERITUM} }\\
    ich & \textbf{trank}\\
    du & \textbf{trankest}\\
    er/sie/es & \textbf{trank}\\
    wir & \textbf{tranken}\\
    ihr & \textbf{trankt}\\
    sie/Sie & \textbf{tranken}\\
  \end{tabular}
  \hfill
  \begin{tabular}{rl}
     \multicolumn{2}{l}{\textbf{IMPERATIV} }\\
   & \textbf{trink(e) (du)!}\\
   & \textbf{trinken wir!}\\
   & \textbf{trinkt (ihr)!}\\
   & \textbf{trinken Sie!}\\
   & \vspace{10pt}
  \end{tabular}

    \vspace{10pt}

 \begin{tabular}{rl}
     \multicolumn{2}{l}{\textbf{PERFEKT}}\\
   & haben + \textbf{getrunken}\\
  \end{tabular}
\hfill
   \begin{tabular}{rl}
   
    \multicolumn{2}{l}{\textbf{FUTURE I} }\\
   & werden + \textbf{trinken}\\
 
  \end{tabular}
\hfill
   \begin{tabular}{rl}
      \multicolumn{2}{l}{\textbf{PLUSQUAMPERFEKT}}\\
   & hatten + \textbf{getrunken}\\
  \end{tabular}
  
  
\vspace{10pt}

\begin{tabular}{lll}
\multicolumn{3}{l}{\textbf{MIT PRÄPOSITIONEN}}\\
& \textbf{---} & \\
\end{tabular}
\hfill
\begin{tabular}{lll}
\multicolumn{3}{l}{\textbf{TRENNBARE VERBEN}}\\
& \textbf{---} & \\
\end{tabular}
\vspace{-5pt}
\end{verbentry}

% treffen (to meet)
\begin{verbentry}{
  style = {acc},
  toc = {treffen (to meet)},
  name = {treffen},
  english = {to meet},
  type = {irregular, + Akkusativ},
  auxiliary = {haben}
}
 
    
     \vspace{5pt}

    \hrule
    
    \vspace{5pt}

  \begin{tabular}{rl}
     \multicolumn{2}{l}{\textbf{PRÄSENS} }\\
    ich & \textbf{treffe}\\
    du & \textbf{triffst}\\
    er/sie/es & \textbf{trifft}\\
    wir & \textbf{treffen}\\
    ihr & \textbf{trefft}\\
    sie/Sie & \textbf{treffen}\\
  \end{tabular}
  \hfill
     \begin{tabular}{rl}
    \multicolumn{2}{l}{\textbf{PRÄTERITUM} }\\
    ich & \textbf{traf}\\
    du & \textbf{trafst}\\
    er/sie/es & \textbf{traf}\\
    wir & \textbf{trafen}\\
    ihr & \textbf{traft}\\
    sie/Sie & \textbf{trafen}\\
  \end{tabular}
  \hfill
  \begin{tabular}{rl}
     \multicolumn{2}{l}{\textbf{IMPERATIV} }\\
   & \textbf{triff (du)!}\\
   & \textbf{treffen wir!}\\
   & \textbf{trefft (ihr)!}\\
   & \textbf{treffen Sie!}\\
   & \vspace{10pt}
  \end{tabular}

    \vspace{10pt}

 \begin{tabular}{rl}
     \multicolumn{2}{l}{\textbf{PERFEKT}}\\
  &  haben + \textbf{getroffen}\\
  \end{tabular}
\hfill
   \begin{tabular}{rl}
   
    \multicolumn{2}{l}{\textbf{FUTURE I} }\\
  &  werden + \textbf{treffen}\\
 
  \end{tabular}
\hfill
   \begin{tabular}{rl}
      \multicolumn{2}{l}{\textbf{PLUSQUAMPERFEKT}}\\
   & hatten + \textbf{getroffen}\\
  \end{tabular}
  
  
\vspace{10pt}

\begin{tabular}{lll}
\multicolumn{3}{l}{\textbf{MIT PRÄPOSITIONEN}}\\
& \textbf{---} & \\
\end{tabular}
\hfill
\begin{tabular}{lll}
\multicolumn{3}{l}{\textbf{TRENNBARE VERBEN}}\\
& \textbf{---} & \\
\end{tabular}
\vspace{-5pt}
\end{verbentry}

% trösten (to comfort)
\begin{verbentry}{
  style = {default},
  toc = {trösten (to comfort)},
  name = {trösten},
  english = {to comfort},
  type = {regular},
  auxiliary = {haben}
}


\vspace{5pt}
\hrule
\vspace{5pt}

\begin{tabular}{rl}
\multicolumn{2}{l}{\textbf{PRÄSENS}}\\
ich & \textbf{tröste}\\
du & \textbf{tröstest}\\
er/sie/es & \textbf{tröstet}\\
wir & \textbf{trösten}\\
ihr & \textbf{tröstet}\\
sie/Sie & \textbf{trösten}\\
\end{tabular}
\hfill
\begin{tabular}{rl}
\multicolumn{2}{l}{\textbf{PRÄTERITUM}}\\
ich & \textbf{tröstete}\\
du & \textbf{tröstetest}\\
er/sie/es & \textbf{tröstete}\\
wir & \textbf{trösteten}\\
ihr & \textbf{tröstetet}\\
sie/Sie & \textbf{trösteten}\\
\end{tabular}
\hfill
\begin{tabular}{rl}
\multicolumn{2}{l}{\textbf{IMPERATIV}}\\
& \textbf{tröste (du)!}\\
& \textbf{trösten wir!}\\
& \textbf{tröstet (ihr)!}\\
& \textbf{trösten Sie!}\\
\end{tabular}

\vspace{10pt}

\begin{tabular}{rl}
\multicolumn{2}{l}{\textbf{PERFEKT}}\\
& haben + \textbf{getröstet}\\
\end{tabular}
\hfill
\begin{tabular}{rl}
\multicolumn{2}{l}{\textbf{FUTURE I}}\\
& werden + \textbf{trösten}\\
\end{tabular}
\hfill
\begin{tabular}{rl}
\multicolumn{2}{l}{\textbf{PLUSQUAMPERFEKT}}\\
& hatten + \textbf{getröstet}\\
\end{tabular}

\vspace{10pt}

\begin{tabular}{lll}
\multicolumn{3}{l}{\textbf{MIT PRÄPOSITIONEN}}\\
& \textbf{trösten mit} + Dat & (to comfort with)\\
& \textbf{trösten über} + Akk & (to console about)\\
\end{tabular}
\hfill
\begin{tabular}{lll}
\multicolumn{3}{l}{\textbf{TRENNBARE VERBEN}}\\
\vspace{20pt}
\end{tabular}
\vspace{-5pt}
\end{verbentry}

% tun (to do)
\begin{verbentry}{
  style = {acc},
  toc = {tun (to do)},
  name = {tun},
  english = {to do},
  type = {irregular, + Akkusativ},
  auxiliary = {haben}
}


\vspace{5pt}
\hrule
\vspace{5pt}

\begin{tabular}{rl}
\multicolumn{2}{l}{\textbf{PRÄSENS}}\\
ich & \textbf{tue}\\
du & \textbf{tust}\\
er/sie/es & \textbf{tut}\\
wir & \textbf{tun}\\
ihr & \textbf{tut}\\
sie/Sie & \textbf{tun}\\
\end{tabular}
\hfill
\begin{tabular}{rl}
\multicolumn{2}{l}{\textbf{PRÄTERITUM}}\\
ich & \textbf{tat}\\
du & \textbf{tatest}\\
er/sie/es & \textbf{tat}\\
wir & \textbf{taten}\\
ihr & \textbf{tatet}\\
sie/Sie & \textbf{taten}\\
\end{tabular}
\hfill
\begin{tabular}{rl}
\multicolumn{2}{l}{\textbf{IMPERATIV}}\\
& \textbf{tu(e) (du)!}\\
& \textbf{tun wir!}\\
& \textbf{tut (ihr)!}\\
& \textbf{tun Sie!}\\
\end{tabular}

\vspace{10pt}

\begin{tabular}{rl}
\multicolumn{2}{l}{\textbf{PERFEKT}}\\
& haben + \textbf{getan}\\
\end{tabular}
\hfill
\begin{tabular}{rl}
\multicolumn{2}{l}{\textbf{FUTURE I}}\\
& werden + \textbf{tun}\\
\end{tabular}
\hfill
\begin{tabular}{rl}
\multicolumn{2}{l}{\textbf{PLUSQUAMPERFEKT}}\\
& hatten + \textbf{getan}\\
\end{tabular}

\vspace{10pt}

\begin{tabular}{lll}
\multicolumn{3}{l}{\textbf{MIT PRÄPOSITIONEN}}\\
& \textbf{tun für} + Akk & (to do for)\\
& \textbf{tun gegen} + Akk & (to do against / about)\\
\vspace{30pt}
\end{tabular}
\hfill
\begin{tabular}{lll}
\multicolumn{3}{l}{\textbf{TRENNBARE VERBEN}}\\
& \textbf{mit$|$tun} & (to participate / be involved)\\
& \textbf{ab$|$tun} & (to dismiss / remove)\\
& \textbf{weg$|$tun} & (to put away / remove)\\
& \textbf{hinzu$|$tun} & (to add)\\
& \textbf{weh$|$tun} & (to hurt / cause pain)\\
\end{tabular}

\vspace{-5pt}
\end{verbentry}
